This chapter introduces CATS,  a novel end-to-end contrastive learning framework for anomaly detection in time series data. CATS addresses the limitations of traditional contrastive learning methods through two key components: temporal similarity and negative data augmentation. By incorporating negative data augmentation, synthetic anomalies are generated, significantly enhancing the model's capability to effectively detect anomalies.
%By leveraging negative data augmentation, synthetic anomalies are introduced, enhancing the model's ability to detect anomalies effectively. 
Additionally, CATS introduces a novel DTW-based temporal loss, enabling efficient time series representation learning by capturing the temporal patterns inherent in time series data.

Empirical evaluations conducted on benchmark datasets and use-case datasets demonstrate the significant improvements achieved by CATS in anomaly detection compared to competing unsupervised models. Importantly, CATS demonstrates superior performance even in challenging scenarios involving data contamination, showcasing its practical applicability to real-world datasets.

In summary, this chapter highlights the significant advancements offered by CATS in the domain of anomaly detection in time series data, establishing it as a versatile and effective framework for addressing diverse real-world challenges.

